\section{Completion and consequences}
\label{sec:completion}

We compare the two rank bounds for the balanced partition of the matrix
layers.  We then transfer the resulting lower bound to every field of
characteristic zero and give an elementary upper bound of the same order
in the exponent.

\begin{proof}[Proof of Theorem~\ref{thm:main}]
Let a syntactically homogeneous $\Sigma\Pi\Sigma\Pi\Sigma$ circuit over
$\C$ of size $S$ compute $f=\IMM_{n,n}$.  Use the partition from
Proposition~\ref{prop:imm}.  For all sufficiently large $n$,
Propositions~\ref{prop:imm} and~\ref{prop:circuit} give absolute
constants $C_1,C_2>0$ such that
\[
 D\exp(-C_1\sqrt n)
 \le R(f)
 \le 2S^2D\exp(C_2\sqrt n)n^{-\sqrt n/100}.
\]
Since $D>0$, cancellation and square roots yield
\[
 S\ge 2^{-1/2}
       \exp\left(-\frac{C_1+C_2}{2}\sqrt n\right)
       n^{\sqrt n/200}.
\]
After increasing the absolute threshold for $n$, the exponential factor
and $2^{-1/2}$ are absorbed by $n^{\sqrt n/400}$.  Therefore
$S\ge n^{\sqrt n/400}$, as asserted.
\end{proof}

The inner products in Proposition~\ref{prop:imm} were used over
$\C$, but its conclusion concerns the rank of a matrix defined over
the integers.  This gives a field-independent consequence.

\begin{corollary}
\label{cor:characteristic-zero}
There is an absolute integer $n_0$ such that, for every field $K$ of
characteristic zero and every $n\ge n_0$, every syntactically homogeneous
$\Sigma\Pi\Sigma\Pi\Sigma$ circuit over $K$ computing $\IMM_{n,n}$
has at least $n^{\sqrt n/400}$ gates.  Bottom fan-in and bottom support
are unrestricted.
\end{corollary}

\begin{proof}
Fix $n$ and the parameters and partition used in
Proposition~\ref{prop:imm}.  Express $F_{\IMM_{n,n}}$ in ordinary
monomial bases of both its domain and codomain.  The coefficients of
$\IMM_{n,n}$ are integers, multiplication has integral matrix entries,
and differentiation has integral matrix entries: a repeated derivative
acts by a product of falling factorials of nonnegative integers.
Consequently this finite matrix has entries in $\Z$.

Every minor is an integer.  Under the natural map $\Z\longrightarrow K$,
an integer is zero precisely when it was zero in $\Z$, because $K$ has
characteristic zero.  The same holds over $\C$.  Thus the largest order
of a nonzero minor, and hence the rank, is the same over $K$ and over
$\C$.  Changing from ordinary monomials to the normalized bases used in
the complex moment argument does not change that complex rank.
Proposition~\ref{prop:imm} therefore supplies its lower bound over $K$.
Proposition~\ref{prop:circuit} holds over every field, with the same
constants.  The preceding comparison proves the corollary with an
absolute threshold independent of $K$.
\end{proof}

A direct block construction gives an upper bound of the form
$n^{\sqrt n+O(1)}$.  The next count includes variable leaves and
permits the computed block entries to be shared.

\begin{proposition}
\label{prop:upper}
Let $K$ be any field and $n\ge2$.  Put
\[
 t=\lceil\sqrt n\rceil,
 \qquad r=\left\lceil\frac nt\right\rceil.
\]
There is a syntactically homogeneous
$\Sigma\Pi\Sigma\Pi\Sigma$ circuit over $K$ computing $\IMM_{n,n}$
whose number of gates, including leaves, is at most
\begin{equation}
\label{eq:upper-gate-bound}
 2n^3+rn^2+rn^{t+1}+n^{r-1}+1
 \le n^{\sqrt n+4}.
\end{equation}
\end{proposition}

\begin{proof}
Partition the $n$ consecutive matrix layers into $r$ nonempty consecutive
blocks, with lengths $\ell_1,\ldots,\ell_r\le t$.  For block $j$, let
$Y^{(j)}$ be the product of its matrices.  An entry $Y^{(j)}_{a,b}$ is a
sum of $n^{\ell_j-1}$ monomials, one for each choice of the internal
indices of that block, and each monomial has degree $\ell_j$.

At the bottom sum layer, use one identity sum gate for each variable,
sharing it among all its uses.  This requires at most $n^3$ leaves and
$n^3$ bottom sum gates.  For every block and each of its $n^2$ entries,
compute the displayed monomials by lower product gates and add them
at a middle sum gate.  Thus block $j$ uses $n^{\ell_j+1}$ lower product
gates and $n^2$ middle sum gates.  When $\ell_j=1$, the product and sum
gates have one input; all five computational layers are still present.

Set $a_0=a_r=1$.  Matrix multiplication now gives
\[
 \IMM_{n,n}
 =\sum_{(a_1,\ldots,a_{r-1})\in[n]^{r-1}}
   \prod_{j=1}^{r}Y^{(j)}_{a_{j-1},a_j}.
\]
For $r=1$ the right-hand side means the single entry $Y^{(1)}_{1,1}$.
Use one upper product gate for each of the $n^{r-1}$ summands, and one
output sum gate.  Each block-entry gate is shared by all summands that
use it.  The upper products have formal degree
$\ell_1+\cdots+\ell_r=n$, and every middle sum adds monomials of its
block's degree.  The bottom sums have degree one, so the whole circuit
is syntactically homogeneous.  Any gates that do not feed the output
may be discarded.

The gate count is at most
\[
 2n^3+rn^2+\sum_{j=1}^{r}n^{\ell_j+1}+n^{r-1}+1,
\]
which proves the first bound in~\eqref{eq:upper-gate-bound}.  For the
second, note that $r\le t\le n$.  If $n\ge3$, then $t\ge2$ and the
first bound is at most
\[
 3n^3+n^{t+2}+n^{t-1}+1
 \le 3n^{t+2}
 \le n^{t+3}
 \le n^{\sqrt n+4}.
\]
Here $3n^3\le n^{t+2}$ and $n^{t-1}+1\le n^{t+2}$ justify the middle
estimate.  For $n=2$, the first bound is $30\le2^5$ and $t=2$, giving
the same conclusion.
\end{proof}
